\subsection{Bonding patterns in the derivative models}

AdNDP was used to compare localized two-center bonding with delocalized
multicenter bonding in the derivative models.\cite{zubarev:adndp:2008,
tkachenko:adndp2:2019}
For triplet \ce{Pt10^{der}}, the alpha and beta density matrices were analyzed
separately. Each spin channel contains 30 Pt-centered 1c elements, with three
on every Pt atom. Nineteen 3c elements occur on the same sets of Pt atoms in
both channels. Pairing these alpha and beta elements gives 19 3c--2e elements
with combined occupation numbers of 1.984--1.994. The alpha channel contains
two additional 3c elements, each with an occupation number of about 0.994.
They are assigned as two 3c--1e elements and account for the two excess alpha
electrons of the triplet.

This spin-resolved pattern agrees qualitatively with the earlier AdNDP analysis of the
lowest-energy high-spin \ce{Pt10} cluster. That structure was described by 30
one-center pairs, 16 3c--2e bonds, four 3c--1e bonds, and four 4c--1e
bonds.\cite{ugartemendia:2022} The exact numbers are different because the two
\ce{Pt10} structures and spin states are different, but both analyses describe
bonding in pure Pt clusters as mainly delocalized and multicenter.

The AdNDP analysis of the \ce{Pt5Ge5^{GM}} provides a complementary
reference. All its valence pairs were assigned to 15 Pt-centered 1c--2e
elements and 20 Pt--Ge 2c--2e bonds.\cite{ugartemendia:2022}
\red{Two alternative decompositions are currently retained for the
\ce{Pt5Ge5} structure obtained through the \ce{Pt10}-to-\ce{Pt-Ge}
derivation (Supporting Table~\ref{tab:tfg_adndp_summary}). Model A contains
10 Ge-centered 1c--2e elements, 21 Pt--Ge 2c--2e bonds, three Ge--Ge 2c--2e
bonds, and one 4c--2e element; model E contains 20 Ge-centered 1c--2e
elements, 14 Pt--Ge 2c--2e bonds, and three Ge--Ge 2c--2e bonds. HAY QUE
ELEGIR UNA DE LAS DOS SOLUCIONES BEFORE SUBMISSION.} Both alternatives
describe a predominantly localized density, unlike the multicenter bonding in
the pure-Pt structure. The \ce{Pt6Ge4} and \ce{Pt4Ge6} derivative
models also favor mainly one- and two-center elements. Ge substitution
therefore changes the mainly multicenter bonding of the pure-Pt structure into
predominantly localized bonding.
Representative multicenter and two-center AdNDP elements are shown in
Figure~\ref{fig:tfg_adndp_comparison}. Additional detailed information is given
in Supporting Table~\ref{tab:tfg_adndp_summary}.

\begin{figure}[H]
\centering
\begin{tikzpicture}[
  x=1cm,
  y=1cm,
  panel/.style={font=\sffamily\bfseries\small, anchor=north west, fill=white,
    fill opacity=0.86, text opacity=1, inner sep=2pt},
  %label/.style={font=\sffamily\bfseries\small, align=center},
  label/.style={font=\sffamily\small, align=center},
  note/.style={font=\sffamily\scriptsize, align=center, text=black!70,
    text width=0.28\linewidth},
  modela/.style={font=\sffamily\bfseries\scriptsize, align=center, text=red!75!black}
]

% Panel sizes and positions can be tuned directly from here.
\def\adndpw{0.305\linewidth}
\def\yimage{0}
\def\ylabel{-2.42}
\def\ynote{-3.15}

\node[inner sep=0pt] (pt10c) at (-4.35,\yimage)
  {\includegraphics[width=\adndpw]{adndp_pt10_3c.png}};
\node[inner sep=0pt] (pt5ge2c) at (0,\yimage)
  {\includegraphics[width=\adndpw]{adndp_pt5ge_2c2e.png}};
\node[inner sep=0pt] (pt5ge4c) at (4.15,\yimage)
  {\includegraphics[width=\adndpw]{adndp_pt5ge_4c2e.png}};

\node[panel] at ([xshift=-1.72cm,yshift=1.82cm]pt10c.center) {(a)};
\node[panel] at ([xshift=-1.72cm,yshift=1.82cm]pt5ge2c.center) {(b)};
\node[panel] at ([xshift=-1.72cm,yshift=1.82cm]pt5ge4c.center) {(c)};

\node[label] at (-4.15,\ylabel) {\ce{Pt10^{der}}};
\node[note] at (-4.15,\ynote) {representative\\3c topology};
\node[label] at (0,\ylabel) {\ce{Pt5Ge5^{der}}};
\node[note] at (0,\ynote) {representative Pt--Ge\\2c--2e element};
\node[label] at (4.15,\ylabel) {\ce{Pt5Ge5^{der}}};
\node[note] at (4.15,\ynote) {residual\\4c--2e element};
\node[modela] at (4.15,-3.72) {Model A only};

\end{tikzpicture}

\caption{Representative AdNDP elements showing the change from multicenter
metallic bonding in \ce{Pt10^{der}} to predominantly localized two-center
bonding in the Ge-containing models: (a) a 3c element of triplet
\ce{Pt10^{der}}, (b) a Pt--Ge 2c--2e element of \ce{Pt5Ge5^{der}}, and (c) the
4c--2e element found in model A of \ce{Pt5Ge5^{der}}. Panel (a) shows the
spatial form of one of the 3c elements identified in the spin-resolved
analysis. Panel (c) is specific to model A and is absent from model E;
its inclusion depends on the AdNDP decomposition selected.}
\label{fig:tfg_adndp_comparison}
\end{figure}

We next turn to the \ce{CO} binding energies of these structures.
